% ===================================================================
%  Глава 1. Анализ предметной области
%  Типы колонок M{} и L{} объявлены в paper.tex
%  \include автоматически добавляет \newpage перед главой
% ===================================================================

\section*{Глава 1. Анализ предметной области}
\addcontentsline{toc}{section}{Глава 1. Анализ предметной области}
\stepcounter{section}

% -------------------------------------------------------------------
\subsection{Обзор существующих программных продуктов по теме работы}
\label{subsec:review}
% -------------------------------------------------------------------

На современном рынке представлен ряд программных решений для
автоматизации бронирования столиков в ресторанах. Для объективной
оценки существующих систем были определены критерии, которым должно
удовлетворять разрабатываемое приложение: онлайн-бронирование
гостем без звонка в заведение; интерактивная схема зала с отображением
статуса столиков; автоматическая проверка пересечений бронирований;
отправка email-уведомлений с PDF-подтверждением; экспорт отчётов о
занятости в форматах CSV и PDF; отсутствие значительных постоянных
расходов; возможность адаптации логики под конкретное заведение.

В соответствии с перечисленными критериями рассмотрены три системы:
iiko, Restoplace и Yclients~\cite{iiko,restoplace,yclients}.

Система \textbf{iiko} --- российская платформа комплексной автоматизации
ресторана (касса, кухня, склад, персонал). Управление схемой зала
доступно администратору, однако онлайн-бронирование гостем не
предусмотрено: резервации вносятся вручную сотрудником. Стоимость
внедрения составляет от 50\,000 рублей, что нецелесообразно для
заведений, которым требуется только онлайн-бронирование.

Сервис \textbf{Restoplace} --- отечественная облачная платформа,
специализирующаяся на бронировании столиков. Гость самостоятельно
выбирает столик на интерактивной схеме зала через виджет на сайте или
во ВКонтакте. Бесплатный тариф ограничен 300 заявками в месяц.
Принципиальным ограничением является отсутствие генерации
PDF-подтверждений и экспорта отчётов --- функций, требуемых в данном
проекте.

Сервис \textbf{Yclients} --- наиболее распространённая в России
CRM-система для сферы услуг (более 26\,000 подключённых заведений).
Предоставляет виджет онлайн-записи, SMS/email-уведомления,
интеграцию с кассами (АТОЛ, Эвотор). Однако система изначально
ориентирована на запись к мастеру, а не на ресторанное бронирование:
нативная схема зала отсутствует, подписка стоит от 3\,000 рублей
в месяц.

Результаты сравнения систем по выделенным критериям приведены
в таблице~\ref{tab:competitors}.

\begin{table}[h!]
    \centering
    \caption{Сравнительный анализ существующих систем по критериям}
    \label{tab:competitors}
    \begin{threeparttable}
        \begin{tabular}{|L{48mm}|M{20mm}|M{21mm}|M{21mm}|M{33mm}|}
            \hline
            \centering\textbf{Критерий} &
            \textbf{iiko} &
            \textbf{Restoplace} &
            \textbf{Yclients} &
            \textbf{Собственная разработка}
            \\ \hline
            Онлайн-бронирование гостем &
            \centering ---\tnote{*} & \centering + & \centering + &
            \centering\arraybackslash +
            \\ \hline
            Интерактивная схема зала &
            \centering + & \centering + & \centering --- &
            \centering\arraybackslash +
            \\ \hline
            Проверка доступности столиков &
            \centering + & \centering + & \centering + &
            \centering\arraybackslash +
            \\ \hline
            Email-уведомления и PDF-подтверждения &
            \centering --- & \centering --- & \centering --- &
            \centering\arraybackslash +
            \\ \hline
            Экспорт отчётов (CSV, PDF) &
            \centering --- & \centering --- & \centering --- &
            \centering\arraybackslash +
            \\ \hline
            Доступность для малого бизнеса &
            \centering --- & \centering +/--- & \centering +/--- &
            \centering\arraybackslash +
            \\ \hline
            Возможность кастомизации &
            \centering --- & \centering --- & \centering --- &
            \centering\arraybackslash +
            \\ \hline
        \end{tabular}
        \begin{tablenotes}
            \footnotesize
            \item[*] Бронирование вносится вручную сотрудником;
            самостоятельная онлайн-резервация гостем не предусмотрена.
        \end{tablenotes}
    \end{threeparttable}
\end{table}

Анализ показал, что ни одна из рассмотренных систем не удовлетворяет
всем критериям одновременно. Ключевыми незакрытыми потребностями
являются генерация PDF-подтверждений и экспорт отчётов о занятости.
Разработка собственного веб-приложения, полностью адаптированного
к требованиям заведения, является обоснованной и целесообразной.

% -------------------------------------------------------------------
\subsection{Анализ программных инструментов разработки
веб-приложений}
\label{subsec:tools}
% -------------------------------------------------------------------

Перед выбором технологического стека проведён обзор инструментов,
актуальных для современной fullstack-разработки, по четырём уровням.

\textbf{Клиентская часть.} Для построения интерфейсов веб-приложений
широко применяются React, Vue.js и Angular~\cite{react_docs}. React ---
библиотека, основанная на компонентном подходе, с наибольшей
экосистемой. Vue.js прост в освоении, но уступает по объёму готовых
библиотек. Angular ориентирован на крупные корпоративные проекты и
требует значительного времени на конфигурацию. Для проекта выбран
\textbf{React}: предоставляет готовые компоненты для интерактивных
форм бронирования (react-calendar, react-toastify).

\textbf{Серверная часть.} Среди актуальных решений выделяются Node.js,
Django, FastAPI и Spring Boot~\cite{django_docs}. Node.js эффективен
при большом числе соединений, но не имеет встроенных аутентификации
и ORM. FastAPI обеспечивает быструю разработку API, однако лишён
административного интерфейса. Spring Boot масштабируем, но требует
длительной конфигурации. Для проекта выбран \textbf{Django}:
предоставляет встроенную аутентификацию, ORM и административную
панель, что напрямую закрывает требования задания.

\textbf{База данных.} Наиболее распространены PostgreSQL, MySQL и
SQLite~\cite{postgres_docs}. SQLite не поддерживает параллельную
запись и непригодна для production. MySQL уступает PostgreSQL по
строгости транзакций. Для проекта выбрана \textbf{PostgreSQL}:
обеспечивает надёжную обработку конкурентных запросов при
одновременном бронировании одного столика, полностью совместима
с Django ORM.

\textbf{Среда развёртывания.} Традиционная ручная настройка сервера
создаёт риск расхождения окружений. Контейнеризация с помощью
\textbf{Docker} и docker-compose устраняет эту проблему: все сервисы
описываются в едином файле и воспроизводятся на любом
хосте~\cite{docker_docs}.













